% Ad Familiares 13.73 --- headnote
% ad-familiares-13-73, undated (range ~54-44 BC), Rome

\textit{Cicero to Q. Marcius Philippus the proconsul,
written from Rome at a date not securely fixed within
the long run of Book 13 recommendation letters
(\textit{Scr.\ Romae anno post ep.\ irxiv}; the
works.yaml entry places the work in the placeholder
year 54 BC, with an outer range to 44 BC). Philippus has
just returned from his province, where he has had L.
Egnatius (in his absence) and L. Oppius (on the spot)
in his care; Cicero opens with a brisk, formal note of
congratulation on the safe homecoming and a promise of
in-person thanks, and then moves straight to the
business of the letter. With Antipater of Derbe ---
the Lycaonian dynast whose city Pompey had attached to
Rome's client network --- Cicero has not only the bond
of \textit{hospitium} but something stronger; he has
heard that Philippus is violently angry with the man,
and that Antipater's sons are in Philippus's power.
Cicero asks, on the strength of their old connection,
that the sons be given to him; he professes neutrality
on the underlying quarrel, and gives Philippus the
explicit out --- if his good name would be touched at
all, the request lapses. The closing is the standard
shape of the book: gratitude promised, and a request
for word back on what is possible.}
